\lysh{34182}{4}

\begin{alist}
\item Το τετράγωνο $\Theta M E \Delta$ έχει πλευρά $x > 0$ (1).
Το εμβαδόν του $E_1 = x^2$.

Το τετράγωνο $H B Z M$ έχει πλευρά $(3 - x) > 0 \Leftrightarrow x < 3$ (2).
Το εμβαδόν του είναι $E_2 = (3 - x)^2$.

Άρα το συνολικό εμβαδόν $E$ ως συνάρτηση του $x$ είναι:

\[E(x) = E_1 + E_2 = x^2 + (3 - x)^2 = 2x^2 - 6x + 9.\]

Από (1) και (2), το πεδίο ορισμού της συνάρτησης $E(x)$ είναι $A = (0, 3)$.

\item Έχουμε:
\begin{align*}
E(x) \ge \dfrac{9}{2} \ &\Leftrightarrow \\
2x^2 - 6x + 9 \ge \dfrac{9}{2} \ &\Leftrightarrow \\
4x^2 - 12x + 18 \ge 9 \ &\Leftrightarrow \\
4x^2 - 12x + 9 \ge 0 \ &\Leftrightarrow
\end{align*}
$(2x - 3)^2 \ge 0$, που ισχύει για κάθε $x \in (0, 3)$.

\item Έχουμε:
\begin{align*}
E(x) &= \dfrac{9}{2} \stackrel{(\beta)}{\Leftrightarrow} \\
(2x - 3)^2 = 0 \ &\Leftrightarrow \\
2x - 3 = 0 \ &\Leftrightarrow \\
x &= \dfrac{3}{2}.
\end{align*}

Το εμβαδόν γίνεται ελάχιστο αν και μόνο αν $x = \dfrac{3}{2}$, δηλαδή αν και μόνο αν το σημείο $E$ είναι μέσο της $\Delta \Gamma$. Στο $\Delta B \Gamma$ τρίγωνο έχουμε:
$E$ μέσο της $\Delta \Gamma$ και $ME // B\Gamma$, οπότε και το $M$ θα είναι μέσο της $B\Delta$.
\end{alist}

